%% This LaTeX-file was created by <johnnyb> Sun Jun 27 22:55:21 1999
%% LyX 1.0 (C) 1995-1999 by Matthias Ettrich and the LyX Team

%% Do not edit this file unless you know what you are doing.
\documentclass[12pt]{article}
\usepackage[T1]{fontenc}
\usepackage{geometry}
\geometry{verbose,tmargin=0in,bmargin=0.5in,lmargin=1in,rmargin=1in,headheight=0.5in,headsep=0.5in,footskip=0.5in}
\usepackage{fancyhdr}
\pagestyle{fancy}
\setcounter{tocdepth}{2}
\usepackage{setspace}
\doublespacing

\makeatletter


%%%%%%%%%%%%%%%%%%%%%%%%%%%%%% LyX specific LaTeX commands.
%\providecommand{\LyX}{L\kern-.1667em\lower.25em\hbox{Y}\kern-.125emX\@}
%% Special footnote code from the package 'stblftnt.sty'
%% Author: Robin Fairbairns -- Last revised Dec 13 1996
%\let\SF@@footnote\footnote
%\def\footnote{\ifx\protect\@typeset@protect
%    \expandafter\SF@@footnote
%  \else
%    \expandafter\SF@gobble@opt
%  \fi
%}
%\expandafter\def\csname SF@gobble@opt \endcsname{\@ifnextchar[%]
%  \SF@gobble@twobracket
%  \@gobble
%}
%\edef\SF@gobble@opt{\noexpand\protect
%  \expandafter\noexpand\csname SF@gobble@opt \endcsname}
%\def\SF@gobble@twobracket[#1]#2{}

\makeatother

\begin{document}


\title{Taking Advantage of Inheritance Through Empirical Factoring }


\author{Jonathan Bartlett}


\date{December 1998}

\maketitle
\tableofcontents{}


\section{Why Object-Orientation is Important}

Object oriented programming is the current revolution in programming techniques.
Object-oriented programming attempts to make extendable, reusable systems. It
does this by creating a dynamic typing system, where new types and typing rules
can be created by the programmer. This allows the programmer to customize his
environment to better suit the problem, encapsulating the underlying structure
and semantics of the programming language to suit the current structure and
semantics of the problem at hand. However, just like any other methodology,
it can be used both poorly and well. This paper intends to give some of the
goals and criteria for the judgment and production of object-oriented systems.


\subsection{The Usefulness of Operational Definitions}

Traditional definitions of objects, generalizations, etc. in history have generally
been definitions that consist of \emph{what attributes} different objects in
the world have. A chair would be defined as a piece of furniture with four legs,
a seat, and a back. However, the problem with defining a chair in terms of its
attributes is that there are many chairs with three legs, which have no backs,
or have similar variations. What makes them still be called a chair?

Operational definitions are the answer to this problem. Operational definitions,
instead of listing the attributes of an object, list the functions of an object.
An operational definition of a chair would be \emph{something a person sits
on}. Therefore, all sorts of shapes and sizes of furniture can be chairs, as
long as a person can sit on it. This allows a person to be able to treat all
chairs similarly without having to worry about how the chair's structure enables
its function. I don't have to figure out how a bean bag chair holds me up, all
I have to know is that I can sit on it.

Similarly, object-oriented programming allows the programmer to publicly describe
objects based on what they can do (their methods) rather than on what they look
like (their implementation). This allows algorithms to operate on all similarly
functioning classes equally well, without having to know anything about their
implementation. This leads to three main advantages to object-oriented programming.

\begin{itemize}
\item Programs do not have to explicitly know the exact types of data they deal with
in order to use them - they must only know the set of methods which they implement
\item Programmers can add additional classes and integrate them with existing programs
without having to change anything in the existing program (except for a recompile
in some languages)\footnotemark{}%

\item Programmers can declare a general set of methods to which other classes must
conform, enabling other methods and objects to use any of these declarations
without knowing the specific type of the object
\item Programs can minimize dependency problems by being dependent on a common, stable
set of classes.\footnotemark{}%

\end{itemize}\addtocounter{footnote}{-1}\footnotetext{
Singer, 117.
}%
\stepcounter{footnote}\footnotetext{
Berard, 76.
}%


\subsection{Inheriting For Reuse}

There are two main reasons for inheritance in object-oriented systems - interface
reuse and implementation reuse. Implementation reuse is also called private
inheritance. It uses another piece of code as its base, but limits and expands
its interface for a new situation. However, since the interface of the object
was not inherited, it cannot be used in the same algorithms as the superclass.\footnote{
Meyer, http://www.eiffel.com/doc/manuals/technology/oosc/inheritance-design/section\_08.html.
} This kind of inheritance does not produce the same kinds of advantages that
interface reuse brings. In fact, many theorists think that it should not be
allowed at all and should only be simulated through aggregation(is-composed-of)
relationships.

Interface (also called facility) reuse is a much more powerful feature of object-oriented
systems. Interface reuse means that all subclasses must implement all of the
interface of the superclass. Data and implementation can be inherited too, but
that is not the major consideration. Interface reuse allows objects and classes
to be used in algorithms that did not know of their existence. This is done
by defining new classes in terms of existing interfaces.\footnote{
Meyer, http://www.eiffel.com/doc/manuals/technology/oosc/inheritance-design/section\_09.html.
} This kind of reuse gives the greatest amount of flexibility and extensibility
to object-oriented systems. Therefore, this paper will focus on interface reuse
and inheritance and not implementation reuse.


\subsection{Creating Inheritance Trees Empirically}

The object-oriented field needs a set of criteria to judge the correctness of
inheritance trees. Codd's rules of relational databases were meant to bring
the best utilization of the relational design model. Codd's rules were made
to put as much of the burden on the DBMS as possible. The rules of object-oriented
programming are made to put as much of the work on the class hierarchy as possible.
These rules will make the programmer make the best use of object-oriented principles,
and will make the best use of the facilities available within object-oriented
programming languages.

These rules only work when applied to existing class definitions. Class definitions
can be reworked to take better advantage of these rules, but that is outside
the scope of this paper. There are many existing papers describing the process
of creating good class definitions, and I will not duplicate their work.\footnote{
http://www.laputan.org/drc/drc.html is a great example of a paper on this subject,
as well as http://www.cse.unsw.edu.au/\~{}timm/pub/subjects/oois96/rules.html.
} However, the inheritance hierarchy should only be created after all classes
and class methods have been defined. As will be shown, inheritance is an empirical
matter which is decided from examining the commonalities within existing classes
and methods. Inheritance is not decided by formulating conceptions beforehand
and trying to impose inheritance on classes.\footnote{
The empirical method of deciding which classes inherit from others will eliminate
what is, according to Rumbaugh, one of the main problems new object oriented
programmers have with inheritance - using inheritance where aggregation is more
appropriate.
}


\section{Rules of Object-Oriented Inheritance}


\subsection{Finding the Greatest Common Factor\footnote{
In this paper, base classes are often referred to as factors, because they function
as a subset of the methods offered by a subclass. The least common denominator
between subclasses refers to the fact that many classes have certain methods
in common, and there should be superclasses defined which take advantage of
the greatest common set of methods between classes.
}}

Inheritance in object-oriented systems can be viewed as a factoring process.
The way to find the greatest common factor between two numbers is to show all
of the prime factors for which each are composed and then find all of the factors
that exist in both sets. The same method is employed when dealing with inheritance.
A superclass is simply a set of methods which is in common between two already-defined
classes. Thus, a superclass could be considered a \emph{factor} of existing
classes. This paper will refer to factors and superclasses synonymously. Most
work in the field of object-oriented programming has dealt with creating subclasses
from superclasses, which is called subclassing. However, this paper deals with
the empirical process of creating superclasses from subclasses, which I will
call factoring.

Object-oriented systems should have all common and similarly functioning methods
of any two classes factored into at least one common superclass. This means
that any two classes which share \emph{any} method(s) should also share a superclass.
For example, if we wrote a program that could read a variety of file systems
(MSDOS, Mac HFS, ISO9660, etc.), we could put each file system in its own class.
Now, each file system would probably have a getFile(filename) method. Therefore,
all of our file systems would need a common superclass that would have a getFile(filename)
method,\footnote{
This class would not have to define any implementation of the getFile method,
only be a placeholder for it. This could be an actual class, or it could be
an abstract class. The only thing is that it needs to have a getFile(filename)
method defined.
} which would probably be called the Filesystem class.\footnote{
There would probably be several methods that were common to all filesystems
that would be implemented here.
}

This factoring of classes into superclasses has one really great benefit - algorithms
can use a variety of different subclasses without having to know which one it
is using. For example, in our file system case, a program could be passed a
Filesystem object (an instance of any of the subclasses), and it could simply
call getFile without caring which Filesystem subclass it was using, or even
having to know which Filesystem subclasses are available at compile time. Therefore,
as more and more Filesystem subclasses are implemented, programs using the generic
Filesystem class would not have to be modified in order to use the new file
systems. This prevents any algorithm from being locked into a specific implementation
of a class and allows it to use any class that can perform the same function.
According to Singer,

\begin{quotation}
A primary benefit of using specialization relationships in analysis is to improve
the clarity of the model . . . It also allows one to easily refer to the common
properties of many types without enumerating them each time . . . It also allows
use of the supertype in the program to represent uses of the subtype and to
be able to add additional subtypes without affecting any of the properties or
behavior of the supertype.\footnote{
Singer, 117.
}
\end{quotation}
The one problem with this is, however, that all of the common methods among
subclasses must accept the same inputs and have the same runtime errors defined.
For example, the getFile method must be able to take a filename of any length
for all of the file systems. This, however, is not a problem if specific errors
are defined in the superclass to be checked, such as file name too long or file
not found. One thing to notice, however, is that there are two types of factors
involved in factoring classes, feature factors and ISA factors.\footnote{
This is terminology I made up in order to better reflect the functions of the
two types of superclasses. Typically, feature factors are called interfaces
or mixin classes, and ISA factors are called base classes.
}


\subsubsection{ISA Factors}

An ISA factor\footnote{
They are called ISA factors because the classes are a type of the superclass.
} is what most people refer to as a superclass. It means that the classes which
have common methods \emph{are a type of} that common object. In our file system
example, the Filesystem class was an ISA factor, because each class \emph{is
a} Filesystem. It is important for any class to only have one ISA factor.\footnote{
This is also called single inheritance. Having one line of inheritance keeps
the namespace from getting cluttered as well as prevent inheritance relationships
which should be aggregation relationships.
} Often times, ISA factors share data and implementation methods as well as the
interface specifications (methods), although this is not required. 


\subsubsection{Feature Factors}

A feature factor is simply a list of operations that many classes have in common.
They are called feature factors because they simply list methods that are available,
and they are usually grouped around a specific function. The factors themselves
only contain the interface, and contain no data or implementation specifics.
For example, there could be a system where many different object types (text,
graphics, etc.) could be sent to the printer with the printMe(printer) method.
This would mean that there could be a feature factor called Printable \emph{}which
contained the method printMe(printer). Any class that wanted to be able to be
printed would have Printable \emph{}as a factor. Also, since feature factors
are simply a list of abilities of an object, a class could implement as many
feature factors as desired.\footnote{
Henderson-Sellers, 247.
} This is similar to multiple inheritance, but it doesn't have the complexity
involved with inheriting code and data. Feature factors are also called \emph{interfaces}
in Java and \emph{Mixin Classes}\footnote{
implemented as abstract base classes
} in systems that support multiple inheritance. In systems that implement late
binding such as Perl and Smalltalk, feature factors don't need to be explicitly
declared, but they should be at least documented for clarification.


\subsubsection{Deciding Between ISA and Feature Factors}

Theoretically, all factorizations could be done with feature factors instead
of ISA factors. However, many applications benefit from a single inheritance
line of code reuse, because it prevents code rewriting. The best way to decide
which factors should be ISA factors is to figure out which factors will share
the most code between subclasses.


\subsubsection{Using Factors Well}

In order to make the best use of feature factors and ISA factors within a program,
method interfaces should specify the highest superclass or factor available
that can get the job done. For example, in the file system situation, if a program
only used behaviors that were available in the Filesystem supertype yet specified
a specific file system like Mac HFS in its interface, it would lose the flexibility
that was gained through factoring. It would be unnecessarily tied to the specific
implementation.


\subsection{Changing States, not Classes}

The next rule of object-oriented programming is that objects should never change
their class type, only their states within a class. For example, it would be
erroneous for me to create a superclass Person, and then two subclasses of person,
Executive and Worker. The problem is that when an object changes its class,
the operations available change as well. So, if a client class were holding
a Worker, and the person suddenly got promoted to an Executive, the assumed
operations available would change, and client classes would no longer be able
to assume that any of the class's operations were available. It would instead
have to continually test for which type it was holding in order to find out
which operations were available. This would nullify all of the advantages of
object-oriented programming. This can be avoided by having things such as age
and occupation be attributes of a class, and not subclasses of them.\footnote{
Johnson and Footer, http://www.laputan.org/drc/drc.html.
} 

Subclasses \emph{must} differ from their superclass by either overrided methods
or by extending its set of methods. Usually both will occur. Changes in attribute
\emph{values} should not create a change in class, only a change in state. This
prevents breaks in the inheritance chain which are unnecessary and could lead
to multiple inheritance and object state changes. For example, if there was
a class Ball which had a method called bounce(), there should not be subclasses
for red and blue balls, because the operation of the ball has not changed. If
we were to later add subclasses to Ball which could each be red or blue, we
would have to resort to multiple inheritance or interfaces (feature factors)
in order to account for all of the variations. This whole mess could be avoided
by simply adding setColor and getColor methods to Ball and leaving the color
as a changeable attribute, not as a part of the type.\footnote{
Rumbaugh, http://www.rational.com/support/techpapers/omt/joop9302.html.
}


\subsection{Access of Non-Constant Data}

All access to non-constant data should be in the form of a method call, not
a direct memory access. This is actually a class design rule, not an inheritance
rule. However, the impact it has on inheritance warranted its inclusion here. 

The reason for this rule is this - if the implementation of the class ever changes
in such a way as to have that value calculated instead of stored, the other
classes which use that class would have to be rewritten. Inheritance also becomes
difficult because subclasses are forced to share implementation with the superclass,
not just the interface. For example, if a subclass calculated a value instead
of storing it, it would break compatibility with the superclass. Data in a class
should be considered the private property of the class, and access should be
explicitly given or restricted through method calls only.\footnote{
Johnson and Footer, http://www.laputan.org/drc/drc.html.
}


\subsection{Unused Data Rule}

This rule is for ISA factors (superclasses which declare both the interface
and implementation). No subclass of an ISA factor should leave data elements
of the factor unused.\footnote{
This does not mean that the subclass has to explicitly refer to all the data
elements, only that it does not render data elements useless through overriding
methods.
} This creates waste of memory by the system, which can readily cause problems
if the subclass is instantiated several times. Also, it can lead to unintended
side-effects later down the road.

Having unused data is is usually a symptom of either of these two problems:

\begin{itemize}
\item The factor is used as an ISA factor when it should only be a feature factor
\item The factor, although it is an ISA factor, declares too many implementation details.
\end{itemize}
These two situations are very similar, and the diagnosis is mostly a judgment
call. The only difference is in the solutions. The solution to the first problem
is this: First, move the common methods to a feature factor. Then, put the superclass
and the subclass in different inheritance hierarchies,\footnote{
Note that now the two classes will no longer have a superclass/subclass relationship.
} but have them both implement the same feature factor. The solution to the second
problem is to create a third class which is also a subclass of the given ISA
factor. The implementation details in the superclass which conflict with the
subclass would be moved to the newly created class. This normally leaves the
superclass as an abstract class.

For example, a system could have classes for a WagedWorker class and an Executive
class.\footnote{
This example is shown to have other problems as well, but it illustrates the
point well.
} Because the WagedWorker and the Executive both have the computePay method,
the Executive class is declared as a subclass of the WagedWorker class. However,
the WagedWorker's pay is calculated based on the hoursWorked and payPerHour
data elements which are stored in the class, while the Executive's pay is calculated
solely on the basis of salary. However, instances of the Executive class will
still contain the hoursWorked and payPerHour data fields, since they are a part
of the WagedWorker superclass, even though they now have no use. As mentioned,
this could be solved in two ways.

\begin{itemize}
\item The WagedWorker and the Executive classes could be put into separate inheritance
lines, and a feature factor, called PaidJob, could be defined which would be
included by both the WagedWorker and the Executive classes
\item Instead of just being a feature factor, the PaidJob class could be an ISA factor
that only has the declaration of the computePay method without the implementation.
Then, the WagedWorker and the Executive classes would each inherit from it.
\end{itemize}

\subsection{Compatibility of Subclass Methods}

As mentioned earlier, in order to be useful, classes which inherit methods from
other classes (either feature factors or ISA factors) must be compatible with
the same types of data and sequences of operations. What this means is that
everywhere a superclass (or factor) is expected, any subclass could be substituted
and it would work properly. This is known as the Liskov substitution principle.
Essentially, this means that in a given situation, all subclasses of a superclass
must perform their operations in a way that is compatible with the subclass.\footnote{
Eliens had a stricter view on this, saying that the operations had to duplicate
the superclass behaviors exactly. However, this is almost impossible to do.
The method proposed by Liskov is more pragmatic, and leaves the designer of
the class responsible for determining the compatibility with the superclass.
} This means that simply having the same method name and parameters may not be
enough. Subclasses must also produce the same types of errors (subclasses cannot
invent their own), and be able to be executed in the same sequence as the superclass.\footnote{
http://c2.com/cgi/wiki?OoDesignPrinciples
}

Subclasses cannot be more restrictive than superclasses classes on parameters
passed to their methods, they can be more restrictive on their returns. Otherwise,
the client objects would have to test for class types before sending messages,
thereby upsetting the whole object-oriented system.\footnote{
Eliens 282 - 283
} Also, they need to be able to follow the same state transition diagrams as
their superclasses (factors). One of the current deficiencies of object-oriented
programming is that state transition diagrams must be enforced by the programmers,
not the language.


\section{Conclusions}


\subsection{The Intended Result - Stable Intermediate Forms}

Although these rules are useful, they can only have their greatest benefit when
classes are designed with this ordering in mind. Programmers need to be looking
for ways in which methods can be shared among the greatest number of classes.
If this is done, after analyzing classes and deciding on the inheritance hierarchy,
there should develop a set of factors which will remain relatively stable throughout
software versions. Many of the common methods of existing classes will also
be common with future classes. This intermediate classes will then become the
stabilizing force of the software development process. According to Berard,
having these stable intermediate classes are what makes an application scalable
through increasingly complex versions.\footnote{
Berard, 76.
}


\subsection{Empiricism over Mysticism}

As shown, the process of creating inheritance hierarchies should not be a mystical
experience. Instead, it should be grounded on solid principles. Many developers
follow this plan intuitively, but it is better to have it spelled out so that
others do not have to rely on their gut feelings to make good programming choices.
These rules should make the greatest use of object-oriented inheritance without
misusing its ideas.


\subsection{With a Grain of Salt}

All the rules and guidelines known to man will not beat good intuition, common
sense, or performance requirements. These guidelines are meant to put more of
the work on the programming language and less on the programmer. Therefore,
if other considerations are more pressing than ease of programming, or if these
rules get in the way of programming, there is nothing sacred about any of these
rules. As always, use judgment. However, I believe these rules will help programmers
to see where and how their object designs could be restructured.

\begin{thebibliography}{1}
\bibitem{1}Berard, Edward V. Essays on Object-Oriented Software Engineering. Englewood
Cliffs, N.J.: Prentice Hall, 1993.
\bibitem{2}Eliens, Anton. Principles of Object-Oriented Software Development. Wokingham,
England: Addison-Wesley Publishing Company, 1995.
\bibitem{2}Henderson-Sellers, Brian. A Book of Object-Oriented Knowledge: Object-Oriented
Analysis, Design, and Implementation: A New Approach to Software Engineering.
New York: Prentice Hall, 1992.
\bibitem{5}Johnson, Ralph E. and Brian Footer. \char`\"{}Designing Reusable Classes\char`\"{}
in Journal of Object Oriented Programming June/July 1988, Volume 1, Number 2,
pgs 22-35. Available online at http://www.laputan.org/drc/drc.html, accessed
December 7, 1998.
\bibitem{5}Meyer, Bertrand, ISE. \emph{Object-Oriented Software Construction, 2nd Edition.}
Prentice Hall, 1997. Chapter 25, ``Using Inheritance,'' Available Online at
http://www.eiffel.com/doc/manuals/technology/oosc/inheritance-design/index.html,
accessed December 7, 1998.
\bibitem{6}Rumbaugh, Dr. James. ``Disinherited! Examples of Misuse of Inheritance.''
http://www.rational.com/support/techpapers/omt/joop9302.html, February, 1993.
Accessed December 7, 1998.
\bibitem{5}Singer, Gilbert L. Object Technology Strategies and Tactics. New York: SIGS
Books and Multimedia, 1996.
\end{thebibliography}
\end{document}
